\chapternoindex{Введение}

Современное машиностроительное проектирование невозможно представить без использования систем автоматизированного проектирования (САПР), таких как КОМПАС-3D. Ключевым преимуществом параметрического моделирования является наличие истории построения — дерева операций, отражающего последовательность создания трёхмерной модели. Однако при работе с импортированными моделями (например, в формате STEP) или моделями, полученными от сторонних разработчиков, эта история часто отсутствует. Инженер получает лишь «немую» геометрию, что существенно затрудняет её редактирование, параметризацию и повторное использование. Задача восстановления истории построения по финальной геометрии, известная как обратный инжиниринг конструктивной истории, является крайне актуальной для промышленности. Она позволяет не только понять логику создания изделия, но и адаптировать его под новые требования, исправить ошибки или использовать в качестве основы для новых разработок.

Однако данная задача является нетривиальной даже с фундаментальной точки зрения. Одна и та же итоговая форма может быть получена различными последовательностями операций, а локальные изменения геометрии могут иметь множество параметрических интерпретаций. Традиционные алгоритмические методы часто оказываются неэффективны из-за высокой вариативности и сложности трёхмерных форм. В последние годы методы машинного обучения, особенно глубокие нейронные сети, показали впечатляющие результаты в анализе трёхмерных данных. Тем не менее, прямое применение существующих подходов для восстановления именно \textit{истории построения} наталкивается на ряд ограничений, связанных с необходимостью предсказывать не только тип операции, но и её точные числовые параметры.

\textbf{Целью} данной выпускной квалификационной работы является разработка программного модуля «НейроКОМПАС-3D», интегрируемого со средой КОМПАС-3D, который на основе анализа геометрии модели позволяет реконструировать последовательность породивших её операций с использованием методов машинного обучения.

Для достижения поставленной цели необходимо решить следующие \textbf{задачи}:
\begin{enumerate}
    \item Провести анализ предметной области, включая форматы представления трёхмерных данных (STEP, STL) и способы описания операций построения (Steps.txt, JSON), а также существующие подходы к анализу 3D-геометрии нейросетевыми методами.
    \item Сформировать требования к набору данных для обучения нейронной сети и подготовить (очистить, разместить) репрезентативный датасет, содержащий примеры моделей с известной историей построения.
    \item Исследовать различные архитектурные подходы к решению задачи, обосновать выбор наиболее перспективного гибридного подхода, сочетающего нейросетевую сегментацию и детерминированные алгоритмы CAD-ядра.
    \item Разработать и обучить нейросетевые модели (классификаторы и сегментаторы) для основных типов операций (выдавливание, вырезание, скругление, фаска) на базе архитектуры MeshCNN.
    \item Спроектировать и реализовать программный модуль, осуществляющий полный цикл инференса: загрузку модели, её предобработку, вызов нейросети, постобработку результатов и подготовку данных для интеграции с КОМПАС-3D.
    \item Провести экспериментальное тестирование разработанного модуля на реальных моделях, оценить точность распознавания и сегментации операций, выявить ограничения и предложить пути дальнейшего развития.
\end{enumerate}

\textbf{Объектом исследования} является процесс восстановления истории построения трёхмерных параметрических моделей. \textbf{Предмет исследования} — методы машинного обучения, применяемые для анализа геометрии и распознавания формообразующих операций в САПР.

В работе использовались следующие \textbf{методы исследования}: анализ научно-технической литературы и документации, системный подход к проектированию программных систем, методы машинного обучения (глубокие нейронные сети, в частности архитектура MeshCNN), методы обработки трёхмерной геометрии, экспериментальное тестирование и сравнительный анализ.

\textbf{Практическая значимость} работы заключается в создании действующего прототипа программного модуля, способного автоматизировать процесс реинжиниринга моделей в среде КОМПАС-3D. Модуль может быть использован для упрощения сложных моделей, восстановления параметров по импортированной геометрии и обучения работе с системой.

\textbf{Структура работы.} Выпускная квалификационная работа состоит из введения, трёх глав, заключения, библиографического списка и приложений. В первой главе проводится анализ предметной области, формализуется задача и обосновывается актуальность исследования. Вторая глава посвящена проектированию архитектуры решения, описанию структуры данных и выбору нейросетевых моделей. В третьей главе описывается программная реализация модуля, приводятся результаты экспериментального обучения и тестирования на реальных деталях, а также обсуждаются пути дальнейшего развития. В приложениях вынесена техническая документация и дополнительные иллюстративные материалы.

